\paragraph*{04.03.04.03. Persönlichkeiten und Arbeitsstile Dritter beurteilen und zum Nutzen des Projekts einsetzen}
\kcilabel{para:04.03.04.03_PersönlichkeitenUndArbeitsstileDritter}{04.03.04.03}
\label{para:04.03.04.03_PersönlichkeitenUndArbeitsstileDritter}

%%% Task
\par Als \gls{wsl} \gls{ricefw} beurteilte ich Persönlichkeiten und Arbeitsstile von Teammitgliedern und Stakeholdern und nutzte diese strategisch. In diesem komplexen Projektumfeld mit divergierenden Interessen und Machtpositionen war es entscheidend, diese Dynamiken zu verstehen und gezielt einzusetzen.

%%% Action
\par Identifizierte Stakeholdergruppen (Betriebsverantwortliche \gls{ais}—\glspl{iea}, \glspl{ieb}, Geschäftsführung \gls{adac}, Projektteams \gls{r3}/\gls{s4}) und analysierte Interessen, Machtpositionen und Einfluss. Nutzte Interviews, Workshops und Analysen. Entwickelte Kommunikations- und Interaktionsstrategie. In \gls{uat} bezog Betriebsverantwortliche frühzeitig ein, adressierte Bedenken. Unterstützte \gls{fb} beim Rahmenwerk für Datenhoheit im Testmandanten (vgl.\autoref{fig:UATDatenDokumentation} und \autoref{fig:TestmandantDatenhoheit}).

%%% Result
\par Verbesserte Zusammenarbeit zwischen Stakeholdern; kanalisierte unstrukturierte Anfragen des \gls{fb} in konsistente Ausführungsreihenfolge. Klare Verantwortlichkeiten minimiert Konflikte. Projektziele erreicht durch strategische Stakeholder-Interaktion.

%%% Reflect
\par Erkenntnis: Verständnis von Persönlichkeiten und Arbeitsstilen ist essentiell für Zusammenarbeit und Projekterfolg. Proaktive, empathische Stakeholder-Interaktion baut Vertrauen und sichert Veränderungsakzeptanz. Zukünftig: Fähigkeiten weiter einsetzen zur effektiven Stakeholder-Dynamik-Verwaltung.